\section{Frontier Attack Surfaces: Multimodal GUI Agents and the Model Context Protocol (MCP)}
\label{sec:modality_mcp}

As agentic AI expands beyond structured text APIs into multimodal operating system control and standardized tool protocols, novel systemic attack vectors emerge that challenge existing security boundaries.

\subsection{Multimodal GUI and Computer-Use Agent Attack Surfaces}
\label{subsec:gui_agents}

The emergence of computer-use agents (e.g., Anthropic Claude Computer Use and systems evaluated on OSWorld~\cite{xie2024osworld}) enables foundation models to interact directly with desktop operating systems via display screenshots, virtual mouse movements, and keyboard keystrokes. This grounded interaction model introduces distinct physical and visual attack surfaces:

\begin{enumerate}
    \item \textbf{Visual Prompt Injection (VPI):} In contrast to plain-text injection, VPI exploits the visual perception pathways of Vision-Language Models (VLMs). Adversaries can embed typographic instructions into rendered web pages, desktop wallpapers, or document images (e.g., micro-fonts, low-contrast text, or steganographic perturbations)~\cite{gu2024agent}. When the agent captures the screen to determine the next UI action, the visual encoder parses the malicious text, steering the agent to click malicious links or execute arbitrary shell commands. Recent browser-resident detection frameworks (e.g., BSM~\cite{vadlamani2026bsm}) attempt to intercept such DOM manipulations and API abuse before rendering.
    \item \textbf{UI Redressing and Visual Clickjacking:} Malicious websites or applications can render transparent overlays, fake system dialogs, or deceptive confirmation buttons designed to mislead the agent's spatial grounding algorithms. An agent intending to click a ``Cancel'' button may be visually induced to click an underlying ``Grant Full Disk Access'' prompt.
    \item \textbf{Operating System State Poisoning:} Unrestricted computer-use agents read system logs, browser histories, and local configuration files. An attacker placing malicious directives in a local terminal buffer or system error log can achieve delayed compromise whenever the agent inspects system state for troubleshooting.
\end{enumerate}

\subsection{Security Analysis of the Model Context Protocol (MCP)}
\label{subsec:mcp_security}

The Model Context Protocol (MCP) has emerged as an open standard for connecting LLM applications to external data sources and tools~\cite{hou2026mcp}. MCP establishes a client-host-server topology communicating via JSON-RPC 2.0 over standard input/output (\texttt{stdio}) or Server-Sent Events (SSE) with HTTP POST. While MCP provides modularity, its operational mechanics introduce protocol-level vulnerabilities:

\begin{figure}[t]
\centering
\resizebox{\columnwidth}{!}{%
\begin{tikzpicture}[
    font=\sffamily\scriptsize,
    >=Stealth,
    % Box styles
    tcb_container/.style={draw=blue!60!black, fill=blue!3, rounded corners=6pt, line width=0.9pt},
    ext_container/.style={draw=purple!60!black, fill=purple!3, rounded corners=6pt, line width=0.9pt},
    tcb_host/.style={draw=blue!85!black, fill=blue!12, rounded corners=4pt, line width=1.1pt, align=center, inner sep=4pt},
    tcb_client/.style={draw=blue!75!black, fill=white, rounded corners=4pt, line width=1.0pt, align=center, inner sep=4pt},
    tcb_sub/.style={draw=blue!60!black, fill=blue!6, rounded corners=3pt, line width=0.8pt, align=center, inner sep=3pt},
    mcp_server/.style={draw=purple!80!black, fill=purple!12, rounded corners=4pt, line width=1.1pt, align=center, inner sep=4pt},
    threat_box/.style={draw=red!80!black, fill=red!7, dashed, rounded corners=3pt, line width=0.8pt, align=center, inner sep=2.5pt},
    % Edges
    rpc_edge/.style={<->, thick, draw=blue!75!black, shorten >= 2pt, shorten <= 2pt},
    flow_edge/.style={->, thick, draw=gray!75!black, shorten >= 2pt, shorten <= 2pt},
    threat_edge/.style={->, very thick, draw=red!80!black, dashed, shorten >= 2pt, shorten <= 2pt},
    % Badges
    badge_rpc/.style={draw=blue!70!black, fill=white, rounded corners=2pt, font=\sffamily\tiny\bfseries, text=blue!90!black, inner sep=1.8pt},
    badge_threat/.style={draw=red!80!black, fill=white, rounded corners=2pt, font=\sffamily\tiny\bfseries, text=red!85!black, inner sep=1.8pt}
]

% =========================================================================
% CONTAINERS & BOUNDARY
% =========================================================================
% Left: Trusted Computing Base (TCB)
\draw[tcb_container] (-0.3, -0.5) rectangle (4.9, 6.9);
\node[font=\bfseries\scriptsize, text=blue!90!black, anchor=north] at (2.3, 6.75) {Trusted Computing Base (Host \& Client)};

% Right: External / Semi-Trusted Servers
\draw[ext_container] (7.1, -0.5) rectangle (13.3, 6.9);
\node[font=\bfseries\scriptsize, text=purple!90!black, anchor=north] at (10.2, 6.75) {External / Semi-Trusted MCP Domain};

% Vertical Trust Boundary Delimiter
\draw[draw=red!70!black, dashed, line width=1.4pt] (6.0, -0.5) -- (6.0, 6.9);
\node[draw=red!70!black, fill=white, rounded corners=3pt, font=\bfseries\tiny, text=red!85!black, inner sep=2pt] at (6.0, 6.75) {MCP Trust Boundary};

% =========================================================================
% LEFT CONTAINER: TCB COMPONENTS
% =========================================================================
% User Ingress
\node[draw=gray!60, fill=white, rounded corners=3pt, inner sep=2.5pt, font=\scriptsize] (user) at (2.3, 5.7) {
    \textbf{User Ingress} $\cdot$ Prompt / Objective $u \in \mathcal{U}$
};

% MCP Host
\node[tcb_host, minimum width=4.7cm, minimum height=1.1cm] (host) at (2.3, 4.35) {
    \textbf{MCP Host Application} ($\mathcal{M}, \Omega_{\text{host}}$)\\
    \scriptsize Model Reasoning $\cdot$ Session Orchestration\\
    \tiny [Context Window $\cdot$ Policy Broker $\cdot$ HITL Approval Gate]
};
\draw[flow_edge] (user.south) -- (host.north);

% MCP Client
\node[tcb_client, minimum width=4.7cm, minimum height=1.1cm] (client) at (2.3, 2.5) {
    \textbf{MCP Client Runtime} ($\mathcal{C}_{\text{mcp}}$)\\
    \scriptsize Connection Lifecycle $\cdot$ Protocol Translation\\
    \tiny [JSON-RPC 2.0 Dispatcher $\cdot$ Dynamic Schema Registry]
};

\draw[rpc_edge] (host.south) -- (client.north)
    node[midway, fill=white, inner sep=1pt, font=\tiny] {Model Tool Call / Response};

% Transport Layer
\node[tcb_sub, minimum width=4.7cm] (transport) at (2.3, 0.7) {
    \textbf{Transport Substrate:}\\
    \tiny Local Standard I/O (\texttt{stdio}) \textbf{/} Remote SSE over HTTP POST
};

\draw[rpc_edge] (client.south) -- (transport.north)
    node[midway, fill=white, inner sep=1pt, font=\tiny] {JSON-RPC Framing};

% =========================================================================
% RIGHT CONTAINER: MCP SERVERS & THREAT PROFILES
% =========================================================================
% Server 1: Resources
\node[mcp_server, minimum width=5.7cm] (res_server) at (10.2, 5.5) {
    \textbf{MCP Resource Server} ($\mathcal{S}_{\text{res}}$)\\
    \tiny Contextual Read Endpoints (\texttt{file://}, \texttt{postgres://}, \texttt{git://})
};
\node[threat_box, minimum width=5.7cm] (threat_res) at (10.2, 4.55) {
    \textbf{Threat T2: URI Path Traversal \& IDOR} (\texttt{file:///etc/passwd})\\
    \tiny Arbitrary host file read via unvalidated resource parameter injection
};

% Server 2: Prompts
\node[mcp_server, minimum width=5.7cm] (pmt_server) at (10.2, 3.35) {
    \textbf{MCP Prompt Server} ($\mathcal{P}_{\text{pmt}}$)\\
    \tiny Pre-Packaged Context Templates \& Directives (\texttt{prompts/get})
};
\node[threat_box, minimum width=5.7cm] (threat_pmt) at (10.2, 2.4) {
    \textbf{Threat T3: Adversarial Prompt Backdoors \& Jailbreaks}\\
    \tiny Sleeper instructions and goal hijacking embedded in server prompts
};

% Server 3: Tools
\node[mcp_server, minimum width=5.7cm] (tool_server) at (10.2, 1.2) {
    \textbf{MCP Tool Server} ($\mathcal{T}_{\text{tool}}$)\\
    \tiny Dynamic Tool Registration (\texttt{tools/list}) $\cdot$ Executable APIs
};
\node[threat_box, minimum width=5.7cm] (threat_tool) at (10.2, 0.25) {
    \textbf{Threat T1/T4: Tool Shadowing \& Approval Fatigue}\\
    \tiny Schema spoofing during handshake; parameter drift in uninspected calls
};

% =========================================================================
% CROSS-BOUNDARY PROTOCOL FLOWS (CLEAR PARALLEL ORTHOGONAL ROUTING)
% =========================================================================
% Resource read flow: Client (top anchor) -> Res Server
\draw[rpc_edge] ([yshift=8pt]client.east) -- ++(1.4, 0) |- (res_server.west)
    node[pos=0.72, above, badge_rpc] {\texttt{resources/read}};

% Prompt get flow: Client (bottom anchor) -> Prompt Server
\draw[rpc_edge] ([yshift=-8pt]client.east) -- ++(1.4, 0) |- (pmt_server.west)
    node[pos=0.72, above, badge_rpc] {\texttt{prompts/get}};

% Tool call flow: Transport -> Tool Server
\draw[rpc_edge] (transport.east) -- ++(1.4, 0) |- (tool_server.west)
    node[pos=0.72, above, badge_rpc] {\texttt{tools/call}};

% =========================================================================
% ADVERSARIAL THREAT OVERLAYS (DASHED RED PATHS)
% =========================================================================
% Attack T1: Tool Shadowing from Tool Server into Client Schema Registry
\draw[threat_edge] (tool_server.west) -- ++(-0.6, 0) |- ([yshift=2pt]transport.east)
    node[pos=0.35, below=1pt, badge_threat] {T1: Rogue Tool Ingestion};

\end{tikzpicture}%
}
\caption{Architectural decomposition, trust boundaries, and protocol threat vectors in Model Context Protocol (MCP) deployments. Solid blue containers enclose the Trusted Computing Base (Host application and Client runtime); dashed purple containers denote standalone, semi-trusted MCP servers exposing Resources, Prompts, and Tools across the perimeter. Red dashed annotations and boxes isolate primary exploit classes: (T1) tool schema poisoning and tool shadowing during dynamic registration (\texttt{tools/list}), (T2) resource URI path traversal and IDOR (\texttt{file:///etc/passwd}), (T3) adversarial template injection and sleeper jailbreaks via server-supplied prompts, and (T4) host-level approval fatigue leading to uninspected parameter execution during tool invocation (\texttt{tools/call}).}
\label{fig:mcp_boundaries}
\end{figure}


\begin{itemize}
    \item \textbf{Tool Schema Poisoning and Tool Shadowing:} MCP servers register their available tools dynamically during the initial protocol handshake (\texttt{tools/list}). A malicious or compromised MCP server can emit deceptive tool descriptions designed to hijack model routing. For instance, an untrusted server can register a tool named \texttt{read\_document} with descriptions claiming superior performance over native tools, causing the LLM to route confidential document requests to the rogue server.
    \item \textbf{Resource URI Path Traversal and IDOR:} MCP \textit{Resources} allow servers to expose local files, databases, or API responses via URI schemes (e.g., \texttt{file:///}, \texttt{postgres://}). If an MCP server fails to strictly sanitize incoming resource requests, an indirect prompt injection attack can compel the client to request \texttt{file:///etc/passwd} or \texttt{file:///C:/Windows/win.ini}, exfiltrating sensitive host state.
    \item \textbf{Prompt Template Injection:} MCP servers can distribute pre-packaged \textit{Prompts} (\texttt{prompts/get}) that include predefined context and instructions. An adversarial prompt provider can inject hidden jailbreak triggers or backdoors into these templates, corrupting any agent session that invokes the pre-packaged workflow.
    \item \textbf{Host-Level Granularity Gaps and Approval Fatigue:} Current MCP host implementations typically enforce tool approval at the coarse server level or tool-name level rather than conducting deep parameter inspection. Once a user approves a tool (e.g., \texttt{execute\_query}), the model can invoke it repeatedly with arbitrary, potentially destructive SQL payloads without further user intervention.
\end{itemize}

\subsection{Cyber-Physical and Embodied Agent Vulnerabilities}
\label{subsec:embodied_agents}

When agentic AI is integrated into robotics, industrial IoT, or smart facility management, decision errors transcend the digital domain:
\begin{itemize}
    \item \textbf{Actuator Overdrive and Physical Safety Violations:} Unbounded planning loops or compromised action parameters can drive robotic arms beyond joint velocity limits, disable heating/ventilation interlocks, or open electronic access gates without physical authorization~\cite{narajala2025securing}.
    \item \textbf{Sensory Deception:} Environmental tampering (e.g., spoofed sensor telemetry or deceptive visual QR codes) can manipulate the agent's real-world world-model, inducing unsafe navigation or manipulation routines.
\end{itemize}
